%%% The main file. It contains definitions of basic parameters and includes all other parts.

%% Settings for single-side (simplex) printing
% Margins: left 40mm, right 25mm, top and bottom 25mm
% (but beware, LaTeX adds 1in implicitly)
\documentclass[12pt,a4paper]{report}
\setlength\textwidth{145mm}
\setlength\textheight{247mm}
\setlength\oddsidemargin{15mm}
\setlength\evensidemargin{15mm}
\setlength\topmargin{0mm}
\setlength\headsep{0mm}
\setlength\headheight{0mm}
% \openright makes the following text appear on a right-hand page
\let\openright=\clearpage

%% Settings for two-sided (duplex) printing
% \documentclass[12pt,a4paper,twoside,openright]{report}
% \setlength\textwidth{145mm}
% \setlength\textheight{247mm}
% \setlength\oddsidemargin{14.2mm}
% \setlength\evensidemargin{0mm}
% \setlength\topmargin{0mm}
% \setlength\headsep{0mm}
% \setlength\headheight{0mm}
% \let\openright=\cleardoublepage

%% Generate PDF/A-2u
\usepackage[a-2u]{pdfx}

%% Character encoding: usually latin2, cp1250 or utf8:
\usepackage[utf8]{inputenc}

%% Prefer Latin Modern fonts
\usepackage{lmodern}

%% Further useful packages (included in most LaTeX distributions)
\usepackage{amsmath}        % extensions for typesetting of math
\usepackage{amsfonts}       % math fonts
\usepackage{amsthm}         % theorems, definitions, etc.
\usepackage{bbding}         % various symbols (squares, asterisks, scissors, ...)
\usepackage{bm}             % boldface symbols (\bm)
\usepackage{graphicx}       % embedding of pictures
\usepackage{fancyvrb}       % improved verbatim environment
\usepackage{natbib}         % citation style AUTHOR (YEAR), or AUTHOR [NUMBER]
\usepackage[nottoc]{tocbibind} % makes sure that bibliography and the lists
			    % of figures/tables are included in the table
			    % of contents
\usepackage{dcolumn}        % improved alignment of table columns
\usepackage{booktabs}       % improved horizontal lines in tables
\usepackage{paralist}       % improved enumerate and itemize
\usepackage{xcolor}         % typesetting in color

%%% Basic information on the thesis

% Thesis title in English (exactly as in the formal assignment)
\def\ThesisTitle{Deep-learning architectures for analysing population neural data}

% Author of the thesis
\def\ThesisAuthor{Petr Houška}

% Year when the thesis is submitted
\def\YearSubmitted{2021}

% Name of the department or institute, where the work was officially assigned
% (according to the Organizational Structure of MFF UK in English,
% or a full name of a department outside MFF)
\def\Department{Department of Software and Computer Science Education}

% Is it a department (katedra), or an institute (ústav)?
\def\DeptType{Department}

% Thesis supervisor: name, surname and titles
\def\Supervisor{Mgr. Ján Antolík, Ph.D.}

% Supervisor's department (again according to Organizational structure of MFF)
\def\SupervisorsDepartment{Department of Software and Computer Science Education}

% Study programme and specialization
\def\StudyProgramme{Computer Science}
\def\StudyBranch{Artificial Intelligence}

% An optional dedication: you can thank whomever you wish (your supervisor,
% consultant, a person who lent the software, etc.)
\def\Dedication{%
First and foremost, I'd like to thank my supervisor Mgr. Ján Antolík, Ph.D. for inviting me to a field that was entirely foreign to me at the beginning, providing thorough and ever-friendly guidance along the (not so short) journey. On a similar note, I want to express my deep gratitude to Prof. Daniel A. Butts for his invaluable and incredibly welcoming consultations regarding both the NDN3 library and the field of computational neuroscience in general.

Although I didn't end up doing my thesis under his supervision, RNDr. Milan Straka, Ph.D. deserves recognition for patiently lending me his valuable time while I was trying to figure out my eventual thesis topic.

Last but not least, I want to acknowledge at least some of the people who were, albeit indirectly, instrumental to me finishing this thesis. Be it through friendly nudging, emotional support, or even just being explicitly busy with their work and thus shaming me out of procrastination. In atmospherically random order, thank you: Václav \& Marcela Houškovi, Petra Millarová, Tomáš Thon, Eliška Kopecká, Ondřej Trhoň, and Adéla Syrová.
}

% Abstract (recommended length around 80-200 words; this is not a copy of your thesis assignment!)
\def\Abstract{%
Accurate models of visual system are key for understanding how our brains process visual information. In recent years, DNNs have been rapidly gaining traction in this domain. However, only few studies attempted to incorporate known anatomical properties of visual system into standard DNN architectures adapted from the general machine learning field, to improve their interpretability and performance on visual data.

In this thesis, we optimize a recent biologically inspired deep learning architecture designed for analysis of population data recorded from mammalian primary visual cortex when presented with natural images as stimuli. We reimplement this prior modeling in existing neuroscience focused deep learning framework NDN3 and assess it in terms of stability and sensitivity to hyperparameters and architecture fine-tuning. We proceed to extend the model with components of various DNN models, analysing novel combinations and techniques from classical computer vision deep learning, comparing their effectiveness against the bio-inspired components. 

We were able to identify modifications that greatly increase the stability of the model while securing moderate improvement in overall performance. Furthermore, we document the importance of small hyperparameters adjustments versus architectural advantages that could facilitate further experiments with the examined architectures. All-new model components were contributed to the open-source NDN3 package. 

Overall, this work grounds previous bio-inspired DNN architectures in the modern NDN3 environment, identifying optimal hyper-parametrization of the model, and thus paving path towards future development of these bio-inspired architectures.
}

% 3 to 5 keywords (recommended), each enclosed in curly braces
\def\Keywords{%
{deep-learning}, {computational neuroscience}, {v1 modeling}, {bio-inspired architectures}, {visual computation}
}

%% The hyperref package for clickable links in PDF and also for storing
%% metadata to PDF (including the table of contents).
%% Most settings are pre-set by the pdfx package.
\hypersetup{unicode}
\hypersetup{breaklinks=true}

% Definitions of macros (see description inside)
%%% This file contains definitions of various useful macros and environments %%%
%%% Please add more macros here instead of cluttering other files with them. %%%

%%% Minor tweaks of style

% These macros employ a little dirty trick to convince LaTeX to typeset
% chapter headings sanely, without lots of empty space above them.
% Feel free to ignore.
\makeatletter
\def\@makechapterhead#1{
  {\parindent \z@ \raggedright \normalfont
   \Huge\bfseries \thechapter. #1
   \par\nobreak
   \vskip 20\p@
}}
\def\@makeschapterhead#1{
  {\parindent \z@ \raggedright \normalfont
   \Huge\bfseries #1
   \par\nobreak
   \vskip 20\p@
}}
\makeatother

% This macro defines a chapter, which is not numbered, but is included
% in the table of contents.
\def\chapwithtoc#1{
\chapter*{#1}
\addcontentsline{toc}{chapter}{#1}
}

% Draw black "slugs" whenever a line overflows, so that we can spot it easily.
\overfullrule=1mm

%%% Macros for definitions, theorems, claims, examples, ... (requires amsthm package)

\theoremstyle{plain}
\newtheorem{thm}{Theorem}
\newtheorem{lemma}[thm]{Lemma}
\newtheorem{claim}[thm]{Claim}

\theoremstyle{plain}
\newtheorem{defn}{Definition}

\theoremstyle{remark}
\newtheorem*{cor}{Corollary}
\newtheorem*{rem}{Remark}
\newtheorem*{example}{Example}

%%% An environment for proofs

\newenvironment{myproof}{
  \par\medskip\noindent
  \textit{Proof}.
}{
\newline
\rightline{$\qedsymbol$}
}

%%% An environment for typesetting of program code and input/output
%%% of programs. (Requires the fancyvrb package -- fancy verbatim.)

\DefineVerbatimEnvironment{code}{Verbatim}{fontsize=\small, frame=single}

%%% The field of all real and natural numbers
\newcommand{\R}{\mathbb{R}}
\newcommand{\N}{\mathbb{N}}

%%% Useful operators for statistics and probability
\DeclareMathOperator{\pr}{\textsf{P}}
\DeclareMathOperator{\E}{\textsf{E}\,}
\DeclareMathOperator{\var}{\textrm{var}}
\DeclareMathOperator{\sd}{\textrm{sd}}

%%% Transposition of a vector/matrix
\newcommand{\T}[1]{#1^\top}

%%% Various math goodies
\newcommand{\goto}{\rightarrow}
\newcommand{\gotop}{\stackrel{P}{\longrightarrow}}
\newcommand{\maon}[1]{o(n^{#1})}
\newcommand{\abs}[1]{\left|{#1}\right|}
\newcommand{\dint}{\int_0^\tau\!\!\int_0^\tau}
\newcommand{\isqr}[1]{\frac{1}{\sqrt{#1}}}

%%% Various table goodies
\newcommand{\pulrad}[1]{\raisebox{1.5ex}[0pt]{#1}}
\newcommand{\mc}[1]{\multicolumn{1}{c}{#1}}

%%% Ref helpers
\newcommand{\refsection}[1]{\textit{\nameref{#1}} (section \ref{#1})}
\newcommand{\refattachment}[1]{\texttt{\nameref{#1}} (attachment \ref{#1})}

\newcommand{\tttnameref}[1]{\texttt{\nameref{#1}}}


% Title page and various mandatory informational pages
\begin{document}
%%% Title page of the thesis and other mandatory pages

%%% Title page of the thesis

\pagestyle{empty}
\hypersetup{pageanchor=false}
\begin{center}

\centerline{\mbox{\includegraphics[width=166mm]{../img/logo-en.pdf}}}

\vspace{-8mm}
\vfill

{\bf\Large MASTER THESIS}

\vfill

{\LARGE\ThesisAuthor}

\vspace{15mm}

{\LARGE\bfseries\ThesisTitle}

\vfill

\Department

\vfill

{
\centerline{\vbox{\halign{\hbox to 0.45\hsize{\hfil #}&\hskip 0.5em\parbox[t]{0.45\hsize}{\raggedright #}\cr
Supervisor of the master thesis:&\Supervisor \cr
\noalign{\vspace{2mm}}
Study programme:&\StudyProgramme \cr
\noalign{\vspace{2mm}}
Study branch:&\StudyBranch \cr
}}}}

\vfill

% Zde doplňte rok
Prague \YearSubmitted

\end{center}

\newpage

%%% Here should be a bound sheet included -- a signed copy of the "master
%%% thesis assignment". This assignment is NOT a part of the electronic
%%% version of the thesis. DO NOT SCAN.

%%% A page with a solemn declaration to the master thesis

\openright
\hypersetup{pageanchor=true}
\pagestyle{plain}
\pagenumbering{roman}
\vglue 0pt plus 1fill

\noindent
I declare that I carried out this master thesis independently, and only with the cited
sources, literature and other professional sources. It has not been used to obtain another
or the same degree.

\medskip\noindent
I understand that my work relates to the rights and obligations under the Act No.~121/2000 Sb.,
the Copyright Act, as amended, in particular the fact that the Charles
University has the right to conclude a license agreement on the use of this
work as a school work pursuant to Section 60 subsection 1 of the Copyright~Act.

\vspace{10mm}

\hbox{\hbox to 0.5\hsize{%
In \hbox to 6em{\dotfill} date \hbox to 6em{\dotfill}
\hss}\hbox to 0.5\hsize{\dotfill\quad}}
\smallskip
\hbox{\hbox to 0.5\hsize{}\hbox to 0.5\hsize{\hfil Author's signature\hfil}}

\vspace{20mm}
\newpage

%%% Dedication

\openright

\noindent
\Dedication

\newpage

%%% Mandatory information page of the thesis

\openright

\vbox to 0.5\vsize{
\setlength\parindent{0mm}
\setlength\parskip{5mm}

Title:
\ThesisTitle

Author:
\ThesisAuthor

\DeptType:
\Department

Supervisor:
\Supervisor, \SupervisorsDepartment

Abstract:
\Abstract

Keywords:
\Keywords

\vss}

\newpage

\openright
\pagestyle{plain}
\pagenumbering{arabic}
\setcounter{page}{1}


%%% A page with automatically generated table of contents of the master thesis

\tableofcontents

%%% Each chapter is kept in a separate file
\chapter*{Introduction}
\addcontentsline{toc}{chapter}{Introduction}

\section*{History \& context}
\addcontentsline{toc}{section}{History \& context}
According to the annals of history, the term \textit{computational neuroscience} was coined by Eric L. Schwartz in 1985 when he organized a conference\footnote{\citep{schwartz1990computational}} to provide a summary of the current status of a field at that time known under a variety of names such as neural modeling, brain theory, and neural networks. The central idea of computational neuroscience, to study the brain through mathematical models, is much older than that, however. First modern traces can be found as early as in 1907 when Louis Lapicque introduced\footnote{\citep{Lapicque1907}} the integrate and fire model. While crude, the model managed to describe observed interactions of a nerve fiber long before the exact mechanisms responsible for generation of neuron action potentials were discovered. Doing so, it greatly contributed to our understanding of the brain and laid a strong foundation for further research.

One of the main approaches in computational neuroscience goes under the name \textit{system identification}: first, through experimentation one obtains a large dataset of input and output pairs of a given neural system - e.g. pairs of images and responses of neurons in primary visual cortex. These are subsequently fitted through machine learning methods with a model, in the hope of identifying the computations that the neural system does to translate the inputs to its outputs. 

Such methods effectively enable us to predict the response of a system to an arbitrary plausible input. Doing so is one of the best ways to test our understanding of such a system. While having a model that predicts the response accurately does not necessarily have to mean our understanding of the biological fundamentals is correct, it is a clear signal it is at least plausible. On the other hand, the opposite, when a model based on our knowledge is not accurate, is proof we should revisit our assumptions. Such an approach is particularly effective in early stages of sensory systems whose primary function is to produce a response encoding to a clearly defined input - sensory stimulus. 

Unsurprisingly, going as far back as to the 1960s, visual neuroscience has been close to traditional image processing and has used its toolkit for computational modeling. The influence was not one-sided. Several ideas, such as convolution layers, inspired by biological observations\footnote{\citep{Lindsay_2020}} have found their way back to classical computer vision. In recent years, the combination of advancements in deep learning vision and higher volume of better data in visual neuroscience, have caused increased focus on deep learning inspired neural models. 

Deep learning inspired models are only a rough approximation of biological reality, not only at the level of single neuron biophysics but also in terms of overall network architecture. For example, deep neural network (DNN) architectures typically do not model interactions between neurons within a single layer or do not account for direct nonlinear interactions between neurons, such as those provided by neuromodulation or attentional mechanisms in the biological brain. 

\section*{Motivation}
\addcontentsline{toc}{section}{Motivation}
Regardless, classic DNNs are still immensely useful. Through the ability to fit their many parameters to real data, and access to ever increasing toolkit of high performance methods borrowed from classical machine learning, they allow for effective means of approximating the stimulus-response function in many neural subsystems. And thanks to their abstraction level, rapid experimentation of higher level concepts is also easier. Much like the integrate and fire model that also abstracted over certain details, they can still inform our understanding of the nature of vision processing. 

However, despite the success of DNNs in neuroscience, having outclassed classical system identification methods in most domains, there remains a substantial gap between the predictions these models give and the actual measured neural responses. This is true even in the most peripheral stages of visual processing such as primary visual cortex. Furthermore, the poor interpretability of the DNN models has been identified as a major challenge limiting their usefulness for understanding neural computation.

To address these issues, a recent model by \cite{antolik}, which will be the focus of the present thesis, explored the idea of incorporating more biologically plausible components into the DNN framework. It showed that such an approach leads to a model with fewer parameters, better interpretability, and outperformed, at the time of writing, other state-of-the-art approaches. 

However, the \citeauthor{antolik} \citeyear{antolik} study also showed shortcomings. The model was very sensitive to random initialization, most hyper-parameters of the model were not thoroughly studied, and more direct one-to-one comparison of the biologically inspired versus classical DNN components was missing. Furthermore, since its publishing, several studies using classical DNN techniques managed to demonstrate slight improvement over the \citeauthor{antolik} data. Finally, the model has been implemented in an ad-hoc way in now discontinued ML framework Theano, which poses a problem for further exploration of the bio-inspired DNNs architecture idea.

\section*{Motivation}
\addcontentsline{toc}{section}{Motivation}
The goal of this thesis is to address the outlined outstanding issues with the \citeauthor{antolik} study, in order to provide a stable, fine-tuned, and well characterized implementation in a modern neuroscience oriented DNN framework to enable future experimentation with this bio-inspired architecture. Following contributions are the goals of this thesis:

\begin{itemize}
    \item \textbf{Assess the \citeauthor{antolik} model (and its variants) in terms of stability and sensitivity to hyperparameter finetuning:} Especially on low quantities of noisy data, which is common in the field of sensory neuroscience, models tend to be relatively unstable with respect to both hyperparameters fine-tuning but also random initialization. We want to quantify the effects and hopefully explore ways to mitigate them to ensure any conclusions drawn are not due to a chance but rather architectural properties. As part of this, we will make use of our whole dataset, comparing various architectures across its three separate subsets.
    \item \textbf{Evaluate novel architectures inspired by classical computer vision:} We want to compare the benefits of biologically restricted techniques with more computationally generic approaches from classical deep computer vision and investigate whether hard constraints could be replaced with well chosen regularization schemes. Furthermore, test the impact of several deep learning techniques that were shown to work on classical computer vision problems.
    \item \textbf{Improve upon \citeauthor{antolik} model:} Decrease the gap between the original model and more recent less biologically constrained architectures that demonstrated improved performance. 
    \item \textbf{Contribute new functionality to the NDN3 library toolbox:} We contribute all tools developed to conduct experiments in this thesis upstream to the NDN3 framework or, where not applicable, publish them separately under open source license. The goal is to enable others to iterate on our findings and evaluate the results in a shared environment.
    \item \textbf{(Re)implement and assess the \citeauthor{antolik} model and several of its variants within the NDN3 framework:} Similar to the previous point, implementing and analysing various models within a shared framework aims to enable rapid prototyping, and generally facilitate further research.
    \item \textbf{Identify opportunities for further research:} Since our set of goals is relatively broad, we will not be able to dive deep into every encountered question. As a result, we want to identify opportunities for more focused further research.
\end{itemize}

This way, the present thesis represents a stepping stone that will accelerate the future long-term research program on bio-inspired DNN architectures for neuroscientific data underway in the \citeauthor{antolik} group.

\section*{Thesis structure}
\addcontentsline{toc}{section}{Thesis structure}
This thesis is divided into 4 parts. First, in \hyperref[ch:1]{chapter 1} we provide the theoretical background. We start with a high level overview of the initial part of the visual processing system. Then, we introduce both computational neuroscience generally and the tools it uses for system identification tasks, as well as provide a brief summary of the deep learning techniques that will be relevant for our experiments. In \hyperref[ch:2]{chapter 2}, we introduce the \citeauthor{antolik} model and other especially relevant work in more detail. Second, in chapters \hyperref[ch:3]{chapter 3} we describe the implementation of the additional methods necessary to realize our architectures and the experiments pipeline. Then, in \hyperref[ch:4]{chapter 4}, we introduce our methodology for both the model exploration and results analysis, and finish with a training regime overview.

\hyperref[ch:5]{Chapter 5} with experiments and their results follows. We start by reimplementing the \citeauthor{antolik} model. In the first section, we analyse it in terms of stability and sensitivity to training hyperparameters, attempting to get the best and most stable version possible. Then we test regularization on the fully connected layers, impact of additional non-linearity, and input scaling. In the second section, we move towards larger architectural changes, testing elements from other state of the art models, traditional deep computer vision, but also drawing inspiration from simpler tools of computational neuroscience. We conduct a comparison between various computationally similar architectures differing only by the explicitness of imposed regularization. In the third section we explored the effects of transfer learning, by training one instance of a model on all of the dataset’s three separate subsets pooled together. Further, we train the best variants of architectures from previous sections and compare their results to assess the universality of their improvements.

Finally, the last \hyperref[ch:6]{chapter} offers a summary of our findings, provides lessons learned, and suggests ample opportunities for further research.
\chapter{Title of the first chapter}

An~example citation: \cite{Andel07}

\section{Title of the first subchapter of the first chapter}

\section{Title of the second subchapter of the first chapter}

\chapter{Title of the second chapter}

\section{Title of the first subchapter of the second chapter}

\section{Title of the second subchapter of the second chapter}


\chapter*{Conclusion}
\addcontentsline{toc}{chapter}{Conclusion}


%%% Bibliography
%%% Bibliography (literature used as a source)
%%%
%%% We employ bibTeX to construct the bibliography. It processes
%%% citations in the text (e.g., the \cite{...} macro) and looks up
%%% relevant entries in the bibliography.bib file.
%%%
%%% The \bibliographystyle command selects, which style will be used
%%% for references from the text. The argument in curly brackets is
%%% the name of the corresponding style file (*.bst). Both styles
%%% mentioned in this template are included in LaTeX distributions.

\bibliographystyle{plainnat}    %% Author (year)
% \bibliographystyle{unsrt}     %% [number]

\renewcommand{\bibname}{Bibliography}

%%% Generate the bibliography. Beware that if you cited no works,
%%% the empty list will be omitted completely.

\bibliography{bibliography}

%%% If case you prefer to write the bibliography manually (without bibTeX),
%%% you can use the following. Please follow the ISO 690 standard and
%%% citation conventions of your field of research.

% \begin{thebibliography}{99}
%
% \bibitem{lamport94}
%   {\sc Lamport,} Leslie.
%   \emph{\LaTeX: A Document Preparation System}.
%   2nd edition.
%   Massachusetts: Addison Wesley, 1994.
%   ISBN 0-201-52983-1.
%
% \end{thebibliography}


%%% Figures used in the thesis (consider if this is needed)
\listoffigures

%%% Tables used in the thesis (consider if this is needed)
%%% In mathematical theses, it could be better to move the list of tables to the beginning of the thesis.
\listoftables

%%% Abbreviations used in the thesis, if any, including their explanation
%%% In mathematical theses, it could be better to move the list of abbreviations to the beginning of the thesis.
\chapwithtoc{List of Abbreviations}

%%% Attachments to the master thesis, if any. Each attachment must be
%%% referred to at least once from the text of the thesis. Attachments
%%% are numbered.
%%%
%%% The printed version should preferably contain attachments, which can be
%%% read (additional tables and charts, supplementary text, examples of
%%% program output, etc.). The electronic version is more suited for attachments
%%% which will likely be used in an electronic form rather than read (program
%%% source code, data files, interactive charts, etc.). Electronic attachments
%%% should be uploaded to SIS and optionally also included in the thesis on a~CD/DVD.
%%% Allowed file formats are specified in provision of the rector no. 72/2017.
\appendix
\chapter{Attachments}

\section{First Attachment}

\openright
\end{document}
